% --------------------------------------------------------------------
%  LEARNING BOXES
% --------------------------------------------------------------------
%  The coloured callouts the chapters use. Every title is a macro, so
%  a book can rename a box without restating its geometry:
%
%      \def\utptitlepitfall{Technical note}   % before the input
%
%  To change how a box LOOKS in one book, use one of the named
%  variations at the end of this section, or \renewtcolorbox after the
%  input. To add a box only one book needs, declare it in that book
%  with \utpnewcallout or \utpnewquietbox.

% Box titles -- override any of these before the input line.
\providecommand{\utptitlekeyidea}{Key idea}
\providecommand{\utptitleworked}{Worked example}
\providecommand{\utptitlepitfall}{Common mistake}
\providecommand{\utptitletryit}{Try it yourself}
\providecommand{\utptitlequickcheck}{Check your understanding}
\providecommand{\utptitlerecap}{Chapter recap}
\providecommand{\utptitledeliverable}{Lab deliverable}
\providecommand{\utptitlelabmode}{How much is scaffolded?}
\providecommand{\utptitleintuition}{Plain idea}
\providecommand{\utptitlefirstread}{Core path}
\providecommand{\utptitleglance}{Chapter at a glance}
\providecommand{\utptitledebug}{Debugging task}
\providecommand{\utptitleoptional}{Beyond the core}
\providecommand{\utptitleverification}{Verification checklist}
\providecommand{\utptitletroubleshoot}{Troubleshooting}
\providecommand{\utptitlesourcebasis}{Source basis}

% Explaining an idea
\newtcolorbox{keyidea}{utpplain,breakable,colback=UTPPaleTeal,colframe=UTPTeal, boxrule=0.8pt,arc=2mm,left=2.5mm,right=2.5mm,top=1.8mm,bottom=1.8mm, fonttitle=\bfseries\color{UTPNavy},title=\utptitlekeyidea}
\newtcolorbox{workedexample}{utpplain,breakable,colback=UTPPaleBlue,colframe=UTPBlue, boxrule=0.8pt,arc=2mm,left=2.5mm,right=2.5mm,top=1.8mm,bottom=1.8mm, fonttitle=\bfseries\color{UTPNavy},title=\utptitleworked}
\newtcolorbox{pitfall}{utpplain,breakable,colback=UTPPaleRed,colframe=UTPRed, boxrule=0.8pt,arc=2mm,left=2.5mm,right=2.5mm,top=1.8mm,bottom=1.8mm, fonttitle=\bfseries\color{UTPNavy},title=\utptitlepitfall}
% Asking the reader to do something
\newtcolorbox{tryit}{utpplain,breakable,colback=white,colframe=UTPTeal, boxrule=0.9pt,arc=2mm,left=2.5mm,right=2.5mm,top=1.8mm,bottom=1.8mm, fonttitle=\bfseries\color{white},colbacktitle=UTPTeal,title=\utptitletryit}
\newtcolorbox{quickcheck}{utpplain,breakable,colback=white,colframe=UTPOrange, boxrule=0.9pt,arc=2mm,left=2.5mm,right=2.5mm,top=1.8mm,bottom=1.8mm, fonttitle=\bfseries\color{white},colbacktitle=UTPOrange,title=\utptitlequickcheck}
\newtcolorbox{chapterrecap}{utpplain,breakable,colback=UTPPaleTeal,colframe=UTPGreen, boxrule=0.9pt,arc=2mm,left=2.5mm,right=2.5mm,top=1.8mm,bottom=1.8mm, fonttitle=\bfseries\color{white},colbacktitle=UTPGreen,title=\utptitlerecap}
% Lab structure
\newtcolorbox{labdeliverable}{utpplain,breakable,colback=white,colframe=UTPGreen, boxrule=0.9pt,arc=2mm,left=2.6mm,right=2.6mm,top=1.9mm,bottom=1.9mm, fonttitle=\bfseries\color{white},colbacktitle=UTPGreen,title=\utptitledeliverable}
\newtcolorbox{labmode}{utpplain,breakable,colback=UTPPaleBlue!52,colframe=UTPBlue, boxrule=0.75pt,arc=2mm,left=2.6mm,right=2.6mm,top=1.7mm,bottom=1.7mm, fonttitle=\bfseries\color{white},colbacktitle=UTPBlue,title=\utptitlelabmode}
% Softer, lighter-weight asides
\newtcolorbox{intuition}{utpplain,breakable,colback=white,colframe=UTPOrange!42, boxrule=0.5pt,arc=2mm,left=2.5mm,right=2.5mm,top=1.5mm,bottom=1.5mm, fonttitle=\bfseries\color{UTPOrange!85!black},title=\utptitleintuition}
\newtcolorbox{firstread}{utpplain,breakable,colback=white,colframe=UTPBlue!42, boxrule=0.5pt,arc=2mm,left=2.5mm,right=2.5mm,top=1.5mm,bottom=1.5mm, fonttitle=\bfseries\color{UTPBlue}, title=\utptitlefirstread}
\newtcolorbox{debugtask}{utpplain,breakable,colback=UTPPaleRed!58,colframe=UTPRed, boxrule=0.85pt,arc=2mm,left=2.6mm,right=2.6mm,top=1.9mm,bottom=1.9mm, fonttitle=\bfseries\color{white},colbacktitle=UTPRed,title=\utptitledebug}
\newtcolorbox{analogy}{utpplain,breakable,colback=white,colframe=UTPOrange!38, boxrule=0.5pt,arc=2mm,left=2.5mm,right=2.5mm,top=1.5mm,bottom=1.5mm, fonttitle=\bfseries\color{UTPOrange!85!black},title=Analogy}
\newtcolorbox{coursechoice}{utpplain,breakable,colback=white,colframe=UTPNavy, boxrule=0.8pt,arc=2mm,left=2.8mm,right=2.8mm,top=2mm,bottom=2mm, fonttitle=\bfseries\color{white},colbacktitle=UTPNavy,title=Course choice}
\newtcolorbox{decisionmemo}{utpplain,breakable,colback=UTPPaleOrange!55,colframe=UTPOrange, boxrule=0.85pt,arc=2mm,left=2.6mm,right=2.6mm,top=1.9mm,bottom=1.9mm, fonttitle=\bfseries\color{white},colbacktitle=UTPOrange,title=Decision memo}
\newtcolorbox{libnote}{utpplain,breakable,colback=white,colframe=UTPMuted!48, boxrule=0.5pt,arc=2mm,left=2.5mm,right=2.5mm,top=1.5mm,bottom=1.5mm, fonttitle=\bfseries\color{UTPNavy},title=Reference note}
\newtcolorbox{miniproject}{utpplain,breakable,colback=UTPPaleTeal!52,colframe=UTPGreen, boxrule=0.9pt,arc=2mm,left=2.6mm,right=2.6mm,top=1.9mm,bottom=1.9mm, fonttitle=\bfseries\color{white},colbacktitle=UTPGreen,title=Mini-project}
% Where the material came from: the lecturer's slides, manual or tutorial
% that a chapter is built on.
\newtcolorbox{sourcebasis}{utpplain,breakable,colback=white,colframe=UTPMuted!55, boxrule=0.5pt,arc=2mm,left=2.5mm,right=2.5mm,top=1.5mm,bottom=1.5mm, fonttitle=\bfseries\color{UTPMuted},title=\utptitlesourcebasis}
% Boxes that take a title argument
\newtcolorbox{modelcard}[1]{utpplain,breakable,colback=white,colframe=UTPBlue, boxrule=0.9pt,arc=2mm,left=2.5mm,right=2.5mm,top=1.8mm,bottom=1.8mm, fonttitle=\bfseries\color{white},colbacktitle=UTPBlue,title=#1}
\newtcolorbox{optionaltopic}{utpplain,breakable,colback=white,colframe=UTPMuted!58, boxrule=0.5pt,arc=2mm,left=2.5mm,right=2.5mm,top=1.5mm,bottom=1.5mm, fonttitle=\bfseries\color{UTPMuted},title=\utptitleoptional}
\newtcolorbox{secondpass}{utpplain,breakable,colback=UTPPaleGray!82,colframe=UTPMuted, boxrule=0.75pt,arc=2mm,left=2.6mm,right=2.6mm,top=1.9mm,bottom=1.9mm, fonttitle=\bfseries\color{white},colbacktitle=UTPMuted,title=Second pass}
\newtcolorbox{labproject}[1]{utpplain,breakable,colback=white,colframe=UTPBlue, boxrule=0.9pt,arc=2mm,left=2.6mm,right=2.6mm,top=1.9mm,bottom=1.9mm, fonttitle=\bfseries\color{white},colbacktitle=UTPBlue,title={#1}}
% Evidence that a lab worked, and what to do when it did not
\newtcolorbox{verification}{utpplain,breakable,colback=UTPPaleTeal!65,colframe=UTPGreen, boxrule=0.8pt,arc=2mm,left=2.6mm,right=2.6mm,top=1.8mm,bottom=1.8mm, fonttitle=\bfseries\color{white},colbacktitle=UTPGreen,title=\utptitleverification}
\newtcolorbox{troubleshoot}{utpplain,breakable,colback=UTPPaleOrange!70,colframe=UTPOrange, boxrule=0.8pt,arc=2mm,left=2.6mm,right=2.6mm,top=1.8mm,bottom=1.8mm, fonttitle=\bfseries\color{white},colbacktitle=UTPOrange,title=\utptitletroubleshoot}

% Named variations a book asks for after the input. Each restyles boxes
% that already exist, so the chapters do not change.
%
%   \utpheavycallouts          a solid colour bar behind the titles of
%                              keyidea, workedexample, intuition and
%                              firstread, with a firmer frame.
%   \utpbreakableframes        let utpfigure and utptable split over a
%                              page break, in a lighter blue (figures.tex).
%   \utplabprojectcolour{...}  the labproject box in another colour, with
%                              a little more room inside.
\newcommand{\utplabprojectcolour}[1]{%
  \renewtcolorbox{labproject}[1]{utpplain,breakable,colback=white,colframe=#1, boxrule=0.9pt,arc=2mm,left=2.7mm,right=2.7mm,top=2mm,bottom=2mm, fonttitle=\bfseries\color{white},colbacktitle=#1,title={##1}}%
}
\newcommand{\utpheavycallouts}{%
  \renewtcolorbox{keyidea}{utpplain,breakable,colback=UTPPaleTeal,colframe=UTPTeal, boxrule=0.85pt,arc=2mm,left=2.5mm,right=2.5mm,top=1.8mm,bottom=1.8mm, fonttitle=\bfseries\color{white},colbacktitle=UTPTeal,title=\utptitlekeyidea}%
  \renewtcolorbox{workedexample}{utpplain,breakable,colback=UTPPaleBlue,colframe=UTPBlue, boxrule=0.8pt,arc=2mm,left=2.5mm,right=2.5mm,top=1.8mm,bottom=1.8mm, fonttitle=\bfseries\color{white},colbacktitle=UTPBlue,title=\utptitleworked}%
  \renewtcolorbox{intuition}{utpplain,breakable,colback=UTPPaleOrange,colframe=UTPOrange, boxrule=0.8pt,arc=2mm,left=2.5mm,right=2.5mm,top=1.8mm,bottom=1.8mm, fonttitle=\bfseries\color{white},colbacktitle=UTPOrange,title=\utptitleintuition}%
  \renewtcolorbox{firstread}{utpplain,breakable,colback=white,colframe=UTPBlue, boxrule=0.9pt,arc=2mm,left=2.5mm,right=2.5mm,top=1.8mm,bottom=1.8mm, fonttitle=\bfseries\color{white},colbacktitle=UTPBlue,title=\utptitlefirstread}%
}

% Declare or restyle a named box (including a shared default). Both constructors take the house
% geometry and let a trailing key list override any of it, so a new box
% is one line and still matches the others:
%
%   \utpnewcallout{toolnote}{white}{UTPBlue}{Tool walkthrough}
%   \utpnewcallout[boxrule=0.9pt]{name}{back}{frame}{Title}
%   \utpnewquietbox{referencebox}{UTPMuted!58}{UTPMuted}{Reference note}
%
% A callout is a tinted box under a solid title bar, for something the
% reader has to act on. A quiet box is a light frame with coloured
% title text, for an aside they can read past. If a second book ever
% wants the same box, move its line into this section.
\newcommand{\utpnewcallout}[5][]{%
  \DeclareTColorBox{#2}{}{utpplain,breakable,colback=#3,colframe=#4, boxrule=0.8pt,arc=2mm,left=2.6mm,right=2.6mm,top=1.8mm,bottom=1.8mm, fonttitle=\bfseries\color{white},colbacktitle=#4,title={#5},#1}%
}
\newcommand{\utpnewquietbox}[5][]{%
  \DeclareTColorBox{#2}{}{utpplain,breakable,colback=white,colframe=#3, boxrule=0.5pt,arc=2mm,left=2.5mm,right=2.5mm,top=1.5mm,bottom=1.5mm, fonttitle=\bfseries\color{#4},title={#5},#1}%
}

% Additional reusable teaching boxes.
\providecommand{\utptitleapplication}{Application walkthrough}
\utpnewquietbox{clarification}{UTPBlue!42}{UTPBlue}{Clarification}
\newtcolorbox{sourcealignment}[1]{utpplain,breakable,colback=white,colframe=UTPMuted!55,boxrule=0.5pt,arc=2mm,left=2.5mm,right=2.5mm,top=1.5mm,bottom=1.5mm,fonttitle=\bfseries\color{UTPMuted},title={#1}}
\utpnewcallout{projectrequirement}{UTPPaleBlue}{UTPBlue}{Project requirement}
\utpnewcallout{recommendedworkflow}{UTPPaleTeal}{UTPTeal}{Recommended workflow}
\utpnewquietbox[colback=UTPPaleBlue!45,boxrule=0.7pt,top=1.8mm,bottom=1.8mm]{conceptmap}{UTPBlue!55}{UTPNavy}{Concept map}
\utpnewquietbox{referencebox}{UTPMuted!58}{UTPMuted}{Reference note}
\utpnewcallout[boxrule=0.9pt,top=1.9mm,bottom=1.9mm]{applicationbox}{white}{UTPBlue}{\utptitleapplication}
